\chapter*{Conferencia XIII}
\addcontentsline{toc}{chapter}{LECTURA I - El timbre y el carácter del violín}
SEÑORES: Al retomar la presentación de las salidas del violín como modificadores del tono, aprovecho la ocasión para reiterar que las salidas son más importantes que todos los demás factores combinados. Por lo tanto, las salidas del violín despiertan un gran interés. "Sin salidas, sin violín" es tan cierto como "sin sol, sin luz". De polo a polo, todos los seres vivos aman la luz del sol. Desde los límites habitables del sur hasta los del norte, la humanidad ama los sonidos dulces. En todo el mundo habitable, el área y la posición de las salidas del violín contribuyen a la felicidad humana. En un área reducida de salidas, el amante estético del violín puede encontrar consuelo. En un área ampliada de salidas, el amante del gran volumen puede encontrar disfrute.
En cuanto a su capacidad para lograr una calidad de sonido excepcional, las salidas de violín son supremas e inigualables. Su poder para suprimir el "ruido" supera al de todos los demás modificadores de tono combinados.
¡Benditos sean!
Señor constructor, al abrir esas salidas, le ruego que use esa afilada hoja con cuidado. La dulce música, que observa atentamente su trabajo, le alabará con dulzura a su cliente. Aún hay quienes le pagarán una onza de oro por cada centavo de madera entre la salida pequeña y la grande.
[Existe una gran clase de usuarios de violín que prefieren la calidad del sonido del violín por encima de la mera cantidad de sonido. Sin dudarlo, confieso pertenecer a dicha clase; pero no me entiendan como si estuviera condenando
190

Cantidad de sonido de violín. Por el contrario, creo que la cantidad de sonido de violín es necesaria en ciertas situaciones; y, en tales situaciones, la cantidad de sonido es tan valiosa como un arma de fuego en el oeste; ambas son un mal necesario para la salvación. Lo que sí condeno es el uso de un sonido potente en todas las situaciones y en todas las ocasiones.
Junto al ruidoso violín, el violín de "gran sonido" resulta ofensivo cuando se emplea en situaciones distintas a aquellas para las que está destinado; es decir, en el gran
auditorios. La ofensa del tono es un asunto serio, ya sea que se deba a un tono excesivo o insuficiente. Mi deseo es contribuir a prevenir el uso del violín en situaciones que despierten desprecio tanto hacia el intérprete como hacia el instrumento. Por lo tanto: llevar un violín viejo, cuyo tono ha decaído considerablemente, a un auditorio grande e intentar forzar su otrora virtuoso sonido ante el oído más lejano expectante no solo es un acto de idiotez inexcusable, sino también una muestra de crueldad despiadada. Posiblemente, la posterior retirada del patrocinio se convierta en un método eficaz para prevenir la continuación de tales espectáculos lamentables.
Los fenómenos más asombrosos relacionados con el sonido del violín se deben al área y la posición de las salidas. Incluso la constancia de estos fenómenos despierta asombro. Hay siete de estos fenómenos, y cuatro de ellos no dependen de la calidad inherente de la madera para existir. Ni la fibra dura, ni la fibra blanda, ni la ausencia total de fibra ejercen influencia alguna sobre esos cuatro fenómenos. La existencia de estos cuatro fenómenos depende
191

totalmente en el aire.
Los siete fenómenos relacionados con las salidas de violín son:
(a) La posición de las salidas puede disminuir o aumentar la potencia del tono.
(b) El aumento del área de salidas incrementa el volumen y disminuye la intensidad del tono.
(do)	La disminución del área de salidas aumenta la intensidad
y disminuye el volumen del tono.	
(d) (e) (f)	La disminución del área de salidas reduce el tono. El aumento del área de salidas aumenta el tono. La disminución del área de salidas reduce el ruido.
calidad del tono.	
(gramo)	Ampliar el área de salidas acentúa el ruido
calidad del tono.
En vista de esta exposición de hechos, no es de extrañar el gran interés que suscitan las salidas del violín. Dado que no se dispone de una explicación satisfactoria para el fenómeno de la elevación y la disminución del tono al aumentar o disminuir el área de las salidas, no se ofrece ninguna; pero, de antemano, agradezco a quien pueda aportar una explicación satisfactoria. (a) La posición de las salidas puede disminuir o aumentar
potencia tonal.
Este hecho es susceptible de demostración concluyente. Así pues: Colocar las salidas en las costillas de la parte central del cuerpo, manteniendo las demás condiciones como de costumbre, reduce la potencia del sonido. La razón de dicha reducción es evidente; las salidas colocadas de esta manera no están en la línea de concentración de las ondas sonoras. Esta conclusión se demuestra cerrando dichas salidas,
y colocando a otros en la posición habitual. 192	El

El notable aumento en la potencia tonal tras dicho cambio de posición es una medida de la concentración de ondas sonoras en ese violín en particular; pero no es una medida de dicha concentración para otro violín con una altura de arqueo de la tapa mayor o menor. Nuevamente, colocar las salidas más cerca de los bordes de la caja de resonancia disminuye la potencia tonal, pero el grado de disminución varía con la altura del arqueo; cuanto más bajo sea el arco, menor será la disminución de la potencia tonal. Es obvio que aumentar la altura del arqueo aumenta la concentración de ondas sonoras en la dirección de la unión central de las tapas; por lo tanto, cuanto más alto sea el arco y más cerca estén las salidas de la unión central, mayor será la potencia tonal. En la práctica, encuentro que una variación de i en la distancia desde la unión central hasta las salidas cambia la potencia tonal cuando la altura del arco es igual o mayor que f; por debajo de f, el cambio en
La potencia tonal es menos marcada. En este sentido, el
El patrón de salidas debe variar en el ancho de la curva en los extremos según el grado de potencia tonal deseado. Por lo tanto, con un ancho de curva disminuido, la salida se puede colocar más cerca de la unión central; un ancho de curva aumentado coloca la salida a una mayor distancia.
desde la unión central. Para la posición de las salidas, no se pueden aplicar reglas estrictas.
(b) El aumento del área de salidas incrementa el volumen y disminuye la intensidad del tono. Mi explicación de este fenómeno comienza con un violín bien hecho, de buen modelo y buena madera, y que no tiene salidas de ningún tipo, pero que tiene otras condiciones como
habitual. 193

En esta condición, las placas permanecen inmóviles y sin sonido bajo una presión vigorosa del arco. Su admirable caja de resonancia no puede funcionar. El aire contenido no le permitirá funcionar. La elasticidad de las moléculas de aire dentro del cuerpo de este violín es una fuerza de resistencia enormemente
mayor que la fuerza en las cuerdas. Solo permitiendo la salida de parte de esta energía resistente se puede excitar esta caja de resonancia. Con el propósito de observar el efecto sobre esta caja de resonancia, permitiré la salida de dicha energía resistente en grados gradualmente crecientes. Si mi razonamiento es correcto, entonces la actividad de la caja de resonancia aumentará a medida que aumente la salida de la fuerza resistente. Con este objetivo en mente, delineo las salidas en su posición habitual y corto el círculo más pequeño en sus extremos superiores. La aplicación del arco ahora produce sonido, pero es un sonido de pequeño volumen, de tono bajo y de marcada intensidad en comparación con el volumen. Es claro de inmediato que pequeño
El volumen de este sonido se debe a la fuga de una pequeña cantidad de energía elástica molecular, pero la razón de la intensidad de este tono no es tan clara, mientras que la razón de su tono bajo me lleva al terreno de la mera suposición. Conozco la suposición de los filósofos sobre este punto, pero la suposición carece de mucha satisfacción. Que las longitudes variables de las columnas de aire confinadas producen sonidos de grados de tono igualmente variables se demuestra ampliamente en los tubos de órgano, las trompas, los silbatos de vapor y todos los instrumentos de viento; pero, después de ampliar estas salidas a la mitad del área habitual, ¿cómo podemos explicarlo?
¿Para el tono notablemente elevado que sigue? 194	Es

Es evidente que el mayor volumen de sonido resultante de dicha ampliación se debe a una mayor liberación de energía elástica molecular; asimismo, que dicha mayor liberación de energía, al disminuir la resistencia, permite una mayor amplitud de la actividad de la caja de resonancia; por lo tanto, una mayor fuerza en el golpe aplicado al aire contenido; por consiguiente, un mayor volumen de tono.
Tras ampliar estas salidas a su área habitual, la aplicación del arco produce un aumento considerable en el volumen del sonido. De nuevo, al ampliar estas salidas más allá de su área habitual, el volumen del sonido aumenta aún más, el tono se eleva, pero la intensidad del sonido disminuye y se acentúa la cualidad ruidosa del sonido.
¿Hasta qué punto puede aumentar el área de salidas?
Extender sin arruinar completamente el valor tonal es una de las demostraciones que no he hecho. Es evidente que tal aumento tiene sus límites prácticos. Sabemos que la ausencia total de paredes confinantes permite la dispersión de las ondas sonoras por igual en todas las direcciones; que dicha dispersión opera para disminuir la intensidad del sonido al mínimo. Es evidente que el área de las salidas puede ser tan grande como para permitir una cantidad de dispersión que disminuya enormemente la intensidad del tono. Es un hecho bien conocido en física que el mayor volumen de sonido se produce al aire libre donde no hay paredes confinantes. En vista de estos hechos, es evidente que disminuir el área de salidas opera para disminuir el volumen del tono, para aumentar la intensidad del tono y para disminuir el ruido. En el trabajo de construcción artística de violines, no conozco ningún factor que posea mayor importancia que
195

Área y posición de las salidas. Si se cuenta con un buen material y un buen modelo, el área y la posición de las salidas deben considerarse para satisfacer los diversos gustos de los violinistas. Así, el constructor que solo busca un gran volumen de sonido colocará salidas de gran área lo más cerca posible de la unión central; y el constructor que busca un menor volumen, mayor intensidad y menor ruido, colocará salidas de área pequeña a mayor distancia de la unión central. En mi opinión, es prudente estar preparado para todas las preferencias tonales; por lo tanto, al regular el sonido en violines destinados a quienes priorizan la calidad sobre el volumen, es necesario comenzar con salidas demasiado pequeñas y aumentar su área solo después de la aplicación del arco. En mi experiencia, es un hecho que no se pueden aplicar reglas estrictas para el área de las salidas. Por lo tanto: si se desea un tono dulce e intenso, la salida de una zona amplia contrarresta ese objetivo. Asimismo, los distintos grados de curvatura generan diferentes trayectorias de las ondas sonoras; por consiguiente, la posición de las salidas debe acercarse o alejarse de la unión central en función de las variaciones en la altura de la curvatura; por lo tanto, no se pueden aplicar reglas estrictas para la posición de las salidas.
En este asunto, como se indicó anteriormente, el constructor de violines debe confiar plenamente en la experiencia y la observación. No hay alternativa. En este asunto, esa cosa colosal llamada "ciencia" es una ayuda.
196

menos como un bebé. En el caso del violín con forma de barril, podemos estar seguros de que las salidas no están en la línea de concentración de ondas sonoras. Esta seguridad se basa en el hecho de que la potencia tonal es invariablemente débil en violines de este modelo. La potencia tonal de este violín redondo es solo la de un bebé; mientras que la potencia tonal de este violín de modelo plano, cada una de sus líneas es una línea de belleza, que simplemente sugiere redondez, hace vibrar cada rincón de la casa con un ritmo alegre y contagioso. Sin lugar a dudas, las salidas de este último violín están en la línea de concentración de ondas sonoras. Creo que una prueba continua podría determinar la distancia precisa de las salidas desde la unión central para cada grado de curvatura. Con la intención de producir la máxima potencia tonal, tales cifras serían valiosas.
Numerosas demostraciones han confirmado que aumentar el área de salidas eleva el tono. Confieso que este fenómeno me parece paradójico. No puedo evitar pensar que ampliar las salidas debería tener el efecto contrario en el tono.
[Aquí hay una demostración de que las ideas preconcebidas pueden llevar a uno al error. A menudo, al levantarnos por la mañana en algún lugar extraño para el cual habíamos formado ideas preconcebidas sobre los diferentes puntos cardinales, el sol parece no salir por el este. Puede parecer que sale por el norte, e incluso por el oeste; y puede negarse obstinadamente a salir por el este mientras permanezcamos en ese punto. Así es como me afecta el fenómeno-
197

el objetivo de elevar el tono mediante la ampliación del área de salidas. Agradezco de antemano cualquier explicación satisfactoria sobre cómo elevar así el tono del violín.
La mínima ampliación necesaria para aumentar el volumen y elevar el tono resulta sorprendente. Basta con eliminar una pequeña porción del borde interior de una salida para modificar estas dos cualidades del sonido.
La filosofía para disminuir o acentuar el "ruido" en el sonido del violín mediante la reducción o el aumento del área de salida no es difícil de explicar. La causa del sonido "ruidoso" del violín, como se ha demostrado anteriormente, se debe a pequeñas áreas de madera más gruesa ubicadas en las partes productoras de sonido de la caja de resonancia; y, debido a esta área de producción de sonido limitada, es evidente que el sonido que produce tiene poca potencia. También es evidente que cualquier agente que actúe para disminuir la potencia del sonido del violín podría aniquilar la débil onda de ruido. Por lo tanto, disminuir el área de salida aumenta la dulzura del sonido del violín; asimismo, cuando existe la causa del "ruido", aumentar el área de salida acentúa dicho sonido ruidoso.
La debilidad de la onda sonora es fácil de demostrar. Saque el violín más ruidoso al aire libre y aplíquele el máximo vigor del arco, mientras el oyente se aleja. A medida que aumenta la distancia, el ruido disminuye, mientras que el sonido musical aún permanece nítido. La posición de las salidas ejerce una influencia considerable en la producción de ese zumbido.
198

El carácter resonante en el sonido del violín se describe acertadamente como "tono interno". Si bien las líneas internas defectuosas de las placas son los principales factores que producen este carácter tonal, la posición de las salidas puede aumentar o disminuir dicho sonido retumbante; pero, en los casos en que este sonido retumbante es marcado, no se puede confiar en la posición de las salidas para lograr una solución completa. De hecho, parece no haber otra solución para los casos graves que construir placas nuevas con líneas internas correctas; es decir, líneas que dirijan de forma continua y precisa el movimiento de las ondas sonoras hacia las salidas en lugar de alejarlo de ellas. Producir la máxima potencia tonal cuando las líneas internas dirigen el movimiento de las ondas sonoras desde las salidas es completamente imposible; además, incluso con dichas líneas absolutamente perfectas, colocar las salidas lejos de las líneas de viaje de las ondas sonoras anula la máxima potencia tonal. En la producción de la máxima potencia tonal, las líneas internas de las placas se convierten en factores prodigiosos; mientras que las líneas externas no ejercen ninguna influencia. Por lo tanto, para lograr la máxima potencia sonora, es necesario primero definir las líneas interiores de las placas y, posteriormente, reducir su espesor trabajando en las líneas exteriores. De este modo, el punto más alto del arco longitudinal, visto desde el interior, puede ubicarse en la posición del puente; así, el movimiento de las ondas sonoras no puede desviarse de las salidas cuando estas se sitúan a la distancia práctica más cercana a la junta central.
(12) La profundidad de las costillas es el último factor en la lista de
199

Modificadores de tono a considerar. El principal efecto de la profundidad de las costillas sobre el tono del violín se manifiesta en la altura del sonido; sin embargo, el efecto de este factor sobre el volumen, la intensidad y el brillo del tono es digno de consideración. Si las demás dimensiones permanecen iguales, la disminución de la profundidad de las costillas reduce el volumen del tono, mientras que la intensidad y el brillo
del tono aumentan. En el experimento sobre esto
factor I he disminuido la profundidad de las costillas de 1 y 1 pulgadas, hacia abajo, en grados de 1-16, a una profundidad igual a 1 pulgada; y desde esta profundidad, he reconstruido dichas costillas, mediante adiciones de 1, a una profundidad de 1 y f pulgadas. Tales cambios en la profundidad, al ser dados a un violín en particular, proporcionan prueba concluyente de que acortar la longitud de las columnas de aire perpendiculares dentro del cuerpo del violín opera para elevar el tono, y viceversa, para bajar el tono. Al disminuir la profundidad de las costillas a 1 pulgada, el tono del violín experimenta cambios notables en su carácter; el volumen del tono disminuye enormemente; el tono aumenta enormemente;
y la intensidad y el brillo del tono aumentan considerablemente. Junto con el área y la posición de las salidas, la profundidad de las costillas es el modificador de tono más potente de la lista. En la profundidad extrema de 1 y k pulgadas, el volumen del tono aumenta considerablemente, mientras que la intensidad y el brillo del tono disminuyen considerablemente, y el tono se reduce considerablemente. En mi observación, la variación de 1 a 32 pulgadas en la profundidad de las costillas opera de manera perceptible para cambiar pmMtiUn estas cualidades tonales. En casos de tono hueco y débil, causado por una disminución excesiva del grosor de la caja de resonancia, he aumentado considerablemente el valor del tono simplemente disminuyendo-
200

la profundidad decreciente de las costillas; y la cantidad de dicha disminución siempre está regida por el grado de debilidad del tono en cada caso. Dicha disminución puede ser igual a 1,32 o 1,16, o 1 pulgada cuando el arqueado es igual, y la profundidad de las costillas es igual a 1 y 1 pulgadas. En este trabajo, el punto es disminuir el volumen del tono, y aumentar la intensidad del tono, el brillo del tono y elevar el tono. Tales efectos sobre el tono se producen con certeza, incluso mientras todas las demás dimensiones permanecen iguales. Es mi observación que el "ruido" desaparece del tono del violín a medida que aumenta la debilidad del tono; y no importa si la debilidad es causada por un grosor demasiado grande o una disminución demasiado grande de la madera de la caja de resonancia, o por un área disminuida de las salidas, o por la aplicación de una sordina al puente. Por lo tanto, el tono hueco y débil nunca es un tono ruidoso. [Hay situaciones en las que el volumen del tono puede ser de mayor valor que la intensidad del tono. Se dan situaciones como aquella en la que la música queda ahogada por el ruido de pasos, conversaciones a gritos, el tintineo de vasos, el estallido de corchos y el tintineo de monedas de plata; situaciones en las que la música se desperdicia en el aire del desierto, me refiero al "aire humeante"; situaciones en las que la única esperanza de estética se centra en la corona nevada de óxido de carbono que corona la elegante goleta. En tales situaciones, se recomienda aumentar el área de las salidas.
Factores externos que actúan para disminuir el violín
La potencia tonal es: (1) La sordina. (2) Las cerdas del arco. 201

La filosofía que subyace a la sordina es el único aspecto de interés que esta posee para el estudioso de las  La sordina demuestra claramente el valor fundamental del puente para el sonido del violín; y esta demostración pone de manifiesto que el puente, al transmitir fuerza a la caja de resonancia, funciona mediante vibración. Dado que la sordina disminuye dicha vibración, resulta útil para demostrar que la rigidez de la madera del puente puede ser tan grande que reduce la potencia del sonido. La sordina también es útil para demostrar que la onda sonora desaparece primero al disminuir la potencia del sonido. Por lo tanto, independientemente del grado de énfasis que un violín le dé al ruido, la aplicación de una sordina al puente lo elimina por completo. En consecuencia, cuando existe un problema de sensibilidad en torno a un violín ruidoso, la sordina se convierte en una aliada. Recuerdos vívidos sugieren que dos aliadas podrían ser mejor que una.
(2) Las cerdas del arco, y la vara misma, pueden disminuir tanto el volumen como la intensidad del tono. Recuerdo un suceso asombroso relacionado con las cerdas del arco desgastadas. Un cierto músico se tomó la molestia de viajar un día para consultarme sobre una pérdida inexplicable en la potencia tonal de su violín. Al presentarme su violín y su arco, comentó: "Solía tener un buen y fuerte tono". Al intentar obtener sonido de su violín, al principio pensé que las cerdas del arco estaban untadas con grasa, porque se deslizaban sobre la
202

Las cuerdas no producían tanto tono como con el uso de la sordina. Tampoco ayudó el aumento de presión. El examen del pelo reveló dos condiciones que disminuían la potencia del sonido. Primero: el diámetro del pelo era mínimo. Segundo: las púas o escamas de las que depende la vibración de las cuerdas estaban completamente desgastadas. Dejando a un lado su arco inservible, usé uno de los míos. Su violín sí tenía un buen tono, fuerte. Al preguntarle, descubrí que había usado mucho las cuerdas Mi de acero. Que alguien con la inteligencia suficiente para leer música no pudiera saber la causa de la pérdida de potencia del sonido en este caso es un problema de peculiaridades humanas para el que no ofrezco explicación.
Sin duda, si esta persona hubiera pagado un buen precio a un buen luthier por un buen violín, ahora lo estaría acusando de fraude. Ante tal provocación, se podría comprender la inquietud de los "viejos maestros".
Hoy en día parece casi necesario afirmar que las cerdas del arco deben ser gruesas y fuertes; que las mejores se obtienen del caballo macho; que, debido a que las púas apuntan en una dirección, para asegurar la misma fuerza en los golpes ascendentes y descendentes, la mitad de dichas púas deben apuntar hacia arriba y la otra mitad hacia abajo; que, cuando esas púas se desgastan por la fricción con las cuerdas y con los grumos de resina, entonces se debe voltear el "hilo"; que, cuando
De nuevo usado, todo el "hebrado" debe ser desechado
203

lejos.
Hay algo de verdad en decir: «Es más difícil encontrar un buen arco que un buen violín». Al acercarme a la vejez, repito ese dicho. Pero en cincuenta años de uso del violín, solo recuerdo haber tenido en mis manos lo que yo llamo un buen arco en dos ocasiones. Durante esos años, he tenido en mis manos muchos buenos violines. Lo que yo llamo un buen arco es aquel que se equilibra a siete pulgadas del talón; que regresa a la posición de reposo con la rapidez del acero templado; que tiene el punto más bajo de la curvatura en su tercio superior y bien arriba hacia la punta; uno que parece abrazar instintivamente las cuerdas.
En la próxima hora se presentará el tema de la máxima uniformidad de la potencia tonal.
